\documentclass[interno]{gpt}

\titolo{Preparazione per incontro con Eggon}
\data{2026-11-10}
\orario{18:00}{20:00}
\luogo{Discord}
\redattore{spagnoli}
\verificatore{corro}
%\assenti{spagnoli}

% \modifica{1.0.0}{2026-11-10}{cescobolla}{-}{Approvazione verbale}
%\modifica{0.1.0}{2026-11-10}{spagnoli}{corro}{Prima stesura}

\componenti{biasini,cescobolla,corro,guglielmo,mazzurana,spagnoli}
\begin{document}

\section{Ordine del giorno}
\begin{enumerate}
  \item Considerazioni e dubbi sui requisiti del progetto.
  \item Rendere più trasparente la rendicontazione delle ore.
  \item Automatizzare il registro delle modifiche.
  \item Implementare uno strumento per il controllo ortografico e grammaticale dei documenti.
\end{enumerate}

\clearpage

\section{Svolgimento}

\subsection{Considerazioni e dubbi sui requisiti del progetto}
Per prepararsi all'incontro di mercoledì 11 novembre 2026 con il proponente, il gruppo ha discusso riguardo ai requisiti di massima del progetto e alle modalità di collaborazione, al fine di poter iniziare a produrre una prima bozza del documento \textit{Analisi dei requisiti}.

\subsection{Rendere più trasparente la rendicontazione delle ore}
Il gruppo ha valutato la possibilità di rendere maggiormente dimostrabile la quantità di ore effettivamente dedicate da ciascun componente per la risoluzione di ogni issue. Invece di inserire un semplice commento con il totale delle ore, si intende indicare l'intervallo temporale esatto in cui si è lavorato per completare la issue.

\subsection{Automatizzare il registro delle modifiche}
Il gruppo ha discusso la possibilità di creare e aggiornare il registro delle modifiche dei documenti tramite l'uso di etichette associate alle pull request, così da evitare inserimenti manuali.

\subsection{Implementazione strumento di controllo ortografia documenti}
Il team ha preso in considerazione l'utilizzo di uno strumento per eseguire il controllo ortografico e grammaticale dei documenti ufficiali.


\section{Decisioni}
\begin{decisioni}
    \decisione{Miglioramenti al sistema di rendicontazione delle ore.}
    \decisione{Creazione di un sistema per automatizzare il registro delle modifiche dei documenti attraverso l'uso delle etichette nelle pull request.}
    \decisione{Test dello strumento per il controllo ortografico e grammaticale dei documenti.}
\end{decisioni}

\section{Azioni da intraprendere}
\begin{azioni}
  \azione[\#94]{Stesura verbale interno della riunione odierna.}{spagnoli}{10/11/2026}
  \azione[\#95]{Automazione versionamento e implementazione LanguageTool.}{guglielmo}{16/11/2026}
  \azione[\#93]{Aggiunta cruscotto nelle norme di progetto.}{mazzurana}{16/11/2026}
\end{azioni}

\section{Prossima riunione}
La prossima riunione interna è prevista per lunedì 16 novembre 2026 alle ore 18:00, con lo scopo di iniziare a definire i requisiti di progetto e i relativi casi d'uso.

\end{document}